\documentclass[10pt, twoside]{article}
\usepackage{be-my-concrete, be-my-geometry}
\usepackage{extra-GT}

\begin{document}

\pagestyle{empty}
\begin{titlepage}
        \vfill

        \shadower\huge\noindent Геометрический подход к гомологиям и когомологиям

        \vfill
\end{titlepage}
\input{misc/abstract.tex}
\newpage

\tableofcontents
\newpage

\setcounter{page}{1}
\pagestyle{fancy}

\section{Клеточные комплексы}
Предметом изучения нашего учебники буду преимущественно \emph{клеточные комплексы} (CW-комплексы), в меньшей степени симплициальные комплесксы. 

\definition{Конечное множество $K$ кусочно-гладких шаров в $\R^n$ называется \emph{клеточным комплексом}, если 
\begin{conditions}
    \item\label{CW1} $|K| \eqdef \bigcup_{\sigma \in K} \sigma = \bigsqcup_{\sigma \in K}\Int \sigma $.
    \item\label{CW2} $\forall \sigma \in K \quad \partial\sigma = \bigcup_{\sigma \in L} \sigma$, где $L \subseteq K$.
\end{conditions}
}
Заметим, что в таком определении пункты \labelcref{CW1,CW2} влекут, что для любых $\sigma, \tau \in K \quad \sigma \cap \tau \in \bigcup K.$ При этом мы не предпологаем, что $\sigma \cap \tau \in K.$

\definition{Подмножество $L \subseteq K$ называется \emph{подкомплексом} клеточного комплекса $K$, если оно само является клеточным комплексом.

Существует еще одно определение, согласующееся с определением топологической пары: $(K, L)$ --- это \emph{клеточная пара}.

\emph{Изоморфизмом} $f\colon (K, L) \to (K', L')$ клеточных пар называется кусочно-гладкий гомеоморфизм $f\colon |K| \to |K'|$, где \begin{inumerate} 
    \item $f(|L|) = |L'|$
    \item $\forall \sigma \in K: f(\sigma) \in K'.$
\end{inumerate}
}
В случае когда $K$ --- симплициальный комплекс, тогда существует \emph{симплициальное отображение} $f_{\Delta}\colon (K, L) \to (K', L')$, т.е. $f(K)$ --- симплициальный комплекс, а $f(L) \subseteq L'$
Под прямым произведение клеточных комплексов будем подразумевать $$K \times L = \{\sigma \times \tau | \sigma \in K , \tau \in L\}.$$



\subsection{Введение в категории}
Для формализации многих конструкций в топологии удобно использовать язык \emph{категорий}. Ввёдем самые базовые понятия этой науки.

\begin{definition}
\emph{Категория} $\C$ имеет следующую структуру:
\begin{conditions}
    \item класса \emph{объектов} $\Ob(\C)$;
    \item для любых объектов $X, Y$ множества \emph{морфизмов} $\Hom_{\C}(X, Y)$\footnote{Часто $\Hom_\C$ опускается до $\Hom$, когда понятно о какой категории идет речь.};
    \item для любых объектов $X, Y, Z$ отображения композиции
    \[
        \Hom(X,Y) \times \Hom(Y,Z) \to \Hom(X,Z),
    \]
    \item для каждого объекта $X$ тождественного морфизма $\id_X \in \Hom(X,X)$.
\end{conditions}
Класс всех морфизмов категории $\C$ будем обозначать $\Mor\C$.
Структура категории $\C$ должна удовлетворять следующим аксиомам:
\begin{conditions}
    \item (ассоциативность) для любых морфизмов $f, g, h$ выполнено
    \[
        h \circ (g \circ f) = (h \circ g) \circ f;
    \]
    \item (единица) для любого морфизма $f : X \to Y$ выполнено
    \[
        f \circ \id_X = f, \quad \id_Y \circ f = f.
    \]
\end{conditions}
\end{definition}

\begin{example}
Категория $\mathbf{Top}$: объекты --- топологические пространства; морфизмы --- непрерывные отображения.
\end{example}

\begin{example}
Категория $\mathbf{Ab}$: объекты --- абелевы группы; морфизмы --- гомоморфизмы групп.
\end{example}

Основная причина использования категорий состоит в том, что многие конструкции (в частности, гомологии) оказываются \emph{функториальными}, то есть согласованными с отображениями пространств.

\begin{definition}
\emph{Функтор} $F \colon \C \to \D$ сопоставляет каждому объекту $X \in \Ob(\C)$ объект $F(X) \in \Ob(\D)$; а каждому морфизму $f : X \to Y$ морфизм $F(f) : F(X) \to F(Y)$, так, что сохраняются композиция и тождественные морфизмы:
\[
F(g \circ f) = F(g)\circ F(f), \quad F(\id_X) = \id_{F(X)}.
\]
\end{definition}

\subsection[Категории Bi и Bs]{Категории $\Bi$ и $\Bs$}
Объектами категории $\Bi$ и $\Bs$ являются пары комплексов $(K, L)$. Морфизмы категории $\Bi$ \emph{пораждаются} изоморфизмами пар и вложениями $(K, L) \hookrightarrow (K', L')$. То есть любой морфизм можно представить как конечную композицию изоморфизмов пар и вложений.

Морфизмы категории $\Bs$ пораждаются теми же изоморфизмами пар и вложениями, а также симплициальными отображениями. 

\section{Подразбиения}
Если $L', L$ --- клеточные комплексы и каждый шар $L'$ содержится в каком-то шаре $L$, а также $|L|$ = $|L'|$, то будем говорить, что $L'$ \emph{подразбивает} $L$ и обозначать это $L \bigtriangleleft L$. В категории $\Bi$ и $\Bs$ дела обстоят немного \quotes{технически лучше}, чем в категории симплициальных комплексов, а именно, если $L \subseteq K$ и $L' \bigtriangleleft L$, тогда есть клеточный комплекс $L' \cup K = \{\sigma| \sigma \in L' \text{ или } \sigma \in K\setminus L\}$.

\subsection{Стягивания}
\begin{definition}
    Пусть $(K_0, K)$ --- клеточная пара. Тогда будем говорить, что $K$ \emph{элементарно стягивается} к $K_0$, если $$K = K_0 \cup B^n \cup B^{n-1},$$
    где $B^n$ --- $n$-мерный шар в $K$, а $B^{n-1} = \partial B^n$ и $B^{n-1}$ содержится только в $B^n$. Элементарное стягивание будем обозначать через $K \searrow^e K_0.$
\end{definition}
\begin{definition}
    Комплекс $K$ \emph{стягивается} к $K_0$, если существует конечная последовательность \[
        K = X_m \searrow^e \cdots \searrow^e X_0 = K_0.
    \]
\end{definition}
\begin{lemma}
    Пусть $(K, L \cup K_0) \subseteq (K, L)$ и $K$ стягивается к $K_0$ причем так, что любой элемент стягивания шара из $L$ проходит вдоль шара из $L_0$\footnote{То есть так, чтобы $L \searrow L_0$}. Тогда сущесвует стянивание $(K, L) \searrow (K_0, L_0)$.
\end{lemma}

Включение $(K_0, L_0) \subseteq (K, L)$ будем называть \emph{расширением}, а также композицию расширения и изоморфизма.

\subsection{Подразбиение в категории частных}
Нам бы вообще хотелось, чтобы все расширения в какой-то категории стали изоморфизмами. Для этого есть целый инструмент --- \emph{локализация}.

Пусть $\B$ обозначает категорию $\Bi$ или $\Bs$. Выделим $\Sigma$ --- класс морфизмов расширений, и посмотрим новую категорию $\B[\Sigma^{-1}]$, объекты категории остаются те же, а морфизмы пораждаются $\Mor(\B)$ и символами $e^{-1}$, где $e \in \Sigma$. Геометрически это означает, что мы можем переходить от комплекса к комплексу, двигаясь как по старым стрелкам из $\Mor(\B)$, так и \quotes{против стрелок} расширений.

\subsubsection{Отношение эквиваленции}
Поскольку одни и те же морфизмы могут теперь записываться по-разному, построим между ними эквиваленцию, которая порождается тремя правилами:
\begin{conditions}
    \item \label{equal1} $\forall f, g \in \Mor(\B): h = fg \implies [h] = [f \circ g];$
    \item \label{equal2} $\forall e \in \Sigma: [e \circ e^{-1}] = [e^{-1} \circ e] = [id];$
    \item \label{equal3} $\forall e_1, e_2 \in \Sigma: [(e_2 e_1)^{-1}] = [e_1^{-1} \circ e_2^{-1}].$
\end{conditions}
\begin{remark}
    Условие \labelcref{equal3} избыточно и следует из условий \labelcref{equal1,equal2}
\end{remark}

\subsection{Универсальное свойство}
Категория частных характеризуется универсальным свойством: здесь оно гарантирует, что $\B[\Sigma^{-1}]$ --- самая свободная категория из тех, где морфизмы $\Sigma$ обратимы.

\begin{proposition}
    Для любого функтора $F \colon \B \to \C$ такого, что $F(e)$ --- изоморфизм для любого $e \in \Sigma$ существует единственный функтор $F' \colon \B[\Sigma^{-1}] \to \C$ такой, что диаграмма
    \begin{equation}
        \begin{tikzcd}
            \B \dar["p"] \rar["F"] & \C \\
            \B[\Sigma^{-1}] \urar["F'", swap]
        \end{tikzcd}
    \end{equation}
    коммутирует, то есть $F = F' \circ p$,  где $p \colon \B \to \B[\Sigma^{-1}]$, $f \mapsto [f]$ --- естественный фунтор локализации.
\end{proposition}

Для простоты дальше мы будем игнорировать пары $(K, L)$, где $L \neq \emptyset$, так как \quotes{общая картина мира} получается за счёт минимальной достройки.


\begin{proposition}
    Любой морфизм в $\Bi[\Sigma^{-1}]$ может быть представлен, как $$[e^{-1}\circ f], \quad e \in \Sigma, f \in \Mor(\Bi).$$
\end{proposition}

\begin{lemma}
    Пусть $e \colon J \to K$ --- расширение и $f \colon J \to L \in \Mor(\Bi)$, а $J_0 = K \cup_{J} L$ --- кодекартов квадрат (pushout или амальгамация).\footnote{$J_0$ получается склейкой $K$ и $L$ по их общей части в $J$.} Тогда существует расширение $e_0 \colon L \to J_0$ и морфизм $f_0 \colon K \to J_0$ такой, что диаграмма
    \begin{equation}
        \begin{tikzcd}
            J \rar["e"] \dar["f"] & K \dar["f_0"] \\
            L \rar["e_0"] & J_0
        \end{tikzcd}
    \end{equation}
    коммутирует, то есть $f_0 \circ e = e_0 \circ f.$
\end{lemma}
\begin{proof}
    Заметим, что $e_0$ и $f_0$ --- естественные вложения в $J_0$. Покажем, что $e_0$ действительно расширение. Поскольку $e\colon J \to K$ --- расширение, тогда $K \searrow J$ стягивается. Если мы приклеем комплекс $L$ к $J$, то \quotes{торчащий} кусок $K$ также стянется. Значит $J_0 \searrow L$, а значит вложение $e_0 \colon L \hookrightarrow J_0$ суть расширение.
\end{proof}
Заметим, что в категории $\Bs$ такое невозможно, для этого посмотрим контрпример. Возьмём $K = I \times I$, где $I$ --- единичный отрезок. Также $J = \{0\} \times I \subset K$ --- левая грань, а $L = \{0\} \times \{0\} \subset J$ --- левая нижняя вершина квадрата. Пусть $f \colon J \to L$ схлопывает отрезок в точку. Чтобы коммутативный квадрат замкнулся $f_0 \colon K \to L$ должен стянуть всю левую грань квадрата в точку, чем сделает эту грань вырожденной.

\subsubsection{Подразделение в категории частных}
Пусть $L'$ подразделение $L$. Посмотрим цилиндр $L \times I$ и заменим его верхнее основание $L \times \{1\}$ на $L'$, получим пространство $(L\times [0, 1)) \cup (L' \times \{1\})$. В такое пространство есть два естественных вложения $i \colon L \hookrightarrow (L\times [0, 1)) \cup (L' \times \{1\})$ и $e \colon L' \hookrightarrow (L\times [0, 1)) \cup (L' \times \{1\})$. Любой цилиндр легко стягивается на своё основание: в нашем случае $(L\times [0, 1)) \cup (L' \times \{1\}) \searrow L'$, таким образом $e$ --- расширение, и в категории $\B[\Sigma^{-1}]$ существует морфизм $[e^{-1} \circ i] \colon L \to L'$, называемый \emph{подразделением} и обозначемый $\bigtriangleright(L, L')$.

\begin{lemma}
    Отображение $\bigtriangleright$ \emph{функториально}, то есть 
    \begin{conditions}
        \item $\bigtriangleright(L, L) = \id_L$;
        \item Если $\bigtriangleright(L', L'')$ и $\bigtriangleright(L, L')$. Тогда $\bigtriangleright(L, L'')$.
    \end{conditions}
\end{lemma}
\begin{proof}
    Во-первых, нужно показать, что если $i_0, i_1 \colon K \hookrightarrow K \times I$ вложения, тогда $[i_0] = [i_1]$. Рассмотрим цилиндр $K \times I$ и подразделим его вставив копию $K$ на уровне $t = \rfrac{1}{2}$. Теперь есть три естественных вложения $i_0', i_1', i_{\rfrac{1}{2}}' \colon K \hookrightarrow K \times I$. Определим отражение $$r \colon K \times I \to K \times I, \quad x \times \{t\} \mapsto x \times \{1-t\}.$$ 
    Заметим, что $r \circ i_{\rfrac 1 2}' = i_{\rfrac 1 2 }'$, а также, что $K \times I \searrow K_{\rfrac 1 2}$. Значит, что $i_{\rfrac 1 2}'$ --- это расширение, тогда в категории $\B[\Sigma^{-1}]$ у $[i_{\rfrac 1 2}']$ есть обратный, тогда умножим равенство на этот обратный:
    \begin{equation}
        \begin{aligned}
            [r] \circ [i_{\rfrac 1 2}'] &= i_{\rfrac 1 2}' \\
            [r] &= [\id_{K\times I}].
        \end{aligned}
    \end{equation}
    При этом $r \times i_0' = i_1'$, а значит 
    \begin{equation}
        [r] \circ [i_0'] = [i_1'] \Longleftrightarrow [i_0'] = [i_1'].
    \end{equation}
    
    Во-вторых, пусть $\Delta^2$ --- 2-симплекс с вершинами $v_0, v_1, v_2$ и противоположным им соотвественно рёбрам $\Delta_0^1, \Delta_1^1, \Delta_2^1.$ Пусть $$(L \times \Delta^2)' \coloneq (L \times \Delta^2) \cup (L' \times \Delta_0^1) \cup (L'' \times \{v_2\}). $$
\end{proof}

\end{document}

